\documentclass[10pt]{article}
\usepackage[left=1in,right=1in,top=3cm,bottom=0.35in,%
includeheadfoot,headheight=99pt,headsep=0.2in]{geometry}
\usepackage{times,graphicx,cite,fancyhdr,amsmath,amssymb,captdef}
\usepackage{marvosym}
\renewcommand{\headrulewidth}{0pt}
%----------------  page style -------------------
\pagestyle{fancy}
\fancyhead{}    %clear the head
\fancyhead[L]{
  \begin{minipage}[t]{1.8in}
    \hspace{0.6in}
    \includegraphics[height=0.6in]{anl.png}
  \end{minipage}
}
\fancyhead[C]{}
\fancyhead[R]{%
  \begin{minipage}[t]{1.7in}
    \begin{flushleft}
      {\scriptsize
        \vspace{-1.1in}
        {\bf Qingteng Zhang, Ph.D., Physicist}\\
        \vspace{0.08in}
        {\bf X-ray Science Division}\\
        Argonne National Laboratory \\
        9700 South Cass Avenue\\
        Argonne, IL 60439 \\
        \vspace{0.08in}
        \Telefon: 630-252-6366 \\
        }
    \end{flushleft}
  \end{minipage}
}
\fancyfoot{}    %clear the foot
\fancyfoot[C]{
  \footnotesize{A U.S. Department of Energy laboratory managed by UChicago
    Argonne, LLC}
}
% ==============================================================================
\setlength{\parskip}{2ex}
\setlength{\parindent}{0pt}

\usepackage[hidelinks]{hyperref}
\usepackage{enumitem}
\setlength{\parskip}{0.4em}
\setlength{\parindent}{0pt}
\raggedbottom
\sloppy   % long unbreakable e-mail addresses have no hyphenation points

\begin{document}

\today

Editors\\
\textit{Nano Letters}\\
American Chemical Society

Dear Editors:

Please consider the enclosed Letter for publication in \textit{Nano Letters}: ``Mesoscale Arrest without Progressive Slowing in a Protein Polymer with Reversible $\beta$-Sheet-like Order.''

\textbf{Corresponding authors.} Qingteng Zhang \mbox{(\href{mailto:qzhang234@anl.gov}{qzhang234@anl.gov})}, Arvind Ramanathan \mbox{(\href{mailto:ramanathana@anl.gov}{ramanathana@anl.gov})}, and Gyorgy Babnigg \mbox{(\href{mailto:gbabnigg@anl.gov}{gbabnigg@anl.gov})}.

\textbf{Significance.} Dynamic arrest in a soft material is usually described as a progressive slowing of motion. In the sequence-defined elastin-like polypeptide (VPAVG)$_{30}$ it proceeds differently. Small-angle X-ray photon correlation spectroscopy (SA-XPCS) at 30~$^\circ$C resolves a two-step correlation decay whose fast-decay fraction falls from $f=0.42$--0.71 to $f=0.01$--0.13 over roughly 47~min, while the fast relaxation time shows no systematic drift and keeps a near-diffusive $Q$ dependence. Material leaves the mobile population without slowing beforehand, and it does so at the largest length scales first: $f$ rises monotonically with $Q$ at every elapsed time, by nine standard deviations or more end to end across 76--167~nm. Arrest therefore proceeds by local conversion rather than bulk stiffening, which implies that a partially arrested protein network still contains fully mobile regions. Temperature-dependent SAXS/WAXS on the same construct resolves a reversible transition to a hierarchical assembly carrying $\beta$-sheet-like peaks at 1.38 and 0.72~\AA$^{-1}$ that vanish on cooling, with seven consecutive 6--34~$^\circ$C cycles recovering the low-temperature profile to within 9\%.

\textbf{Novelty.} Prior work on gel-forming ELPs established arrested phase separation and network morphology from time-averaged scattering and rheology. The dynamical signature reported here, a diffusive fast mode whose time scale plateaus while its amplitude declines, has also been observed in a thermoreversible nanoparticle gel [D. Bahadur \textit{et al.}, \textit{J. Chem. Phys.} \textbf{151}, 104902 (2019)], and the manuscript cites that agreement rather than claiming the phenomenology as new. What is new is the length-scale decomposition. Fitting the fast fraction separately in each $Q$ bin shows that arrest is ordered by length scale, reaching 167~nm before 76~nm at every elapsed time, and follows the transfer to near-completion at $f=0.01$. That decomposition requires microsecond framing: with the fast relaxation time between 0.1 and 1.6~ms, a kilohertz acquisition returns only the slow component. We are equally explicit about the limits. The structural and dynamical observations come from separate experiments, so no causal link is claimed. The $\beta$-sheet assignment rests on WAXS peak positions rather than spectroscopy. The 1.6~s maximum delay bounds a fraction arrested on that scale, not an infinite-time nonergodicity factor. The isothermal series comes from one loaded sample. All four are stated in the manuscript, and the ordered contacts are called plausible rather than established junctions.

\textbf{Impact and urgency.} The work spans sequence-defined nanomaterials, protein self-assembly, nonequilibrium soft matter, and coherent X-ray characterization, and gives a nanoscale structure--dynamics framework for responsive protein materials whose $\beta$-sheet-like order switches thermally without permanent aggregation. Because the behavior is sequence-encoded, the strategy transfers to other programmable protein polymers. The result rests on a capability that is not widely available. The XSPA-500k at 8-ID-I frames at 50~kHz, among the fastest in routine operation at a storage ring. The fast relaxation reported here, 0.1 to 1.6~ms, would fall at or below the shortest delay of a conventional kilohertz acquisition, and the amplitude transfer that is the paper's central result would not be visible at all. High-speed detectors and upgraded storage rings have only recently made this practical on weakly scattering biological samples, so a concise, carefully bounded demonstration is timely. The reduced data and all analysis and figure-generation code are already public on GitHub (\href{https://github.com/qzhang234/PubCode-2021-IDPP}{qzhang234/PubCode-2021-IDPP}).

The manuscript is original, is not under consideration elsewhere, and has been approved by all authors. The authors declare no competing financial interest; AI-tool use is disclosed in the Acknowledgments. Thank you for your consideration.

Sincerely,

\vspace{0.1in}
\nopagebreak
Qingteng Zhang, Ph.D., on behalf of all authors

\newpage
\section*{Suggested Reviewers}
\small
The following researchers have relevant expertise in elastin-like polypeptides,
$\beta$-sheet-forming protein and peptide assemblies, small-angle X-ray
scattering, and X-ray photon correlation spectroscopy.

\textbf{Bradley D. Olsen}\\
Alexander and I. Michael Kasser (1960) Professor of Chemical Engineering, Massachusetts Institute of Technology\\
\href{mailto:bdolsen@mit.edu}{bdolsen@mit.edu}\\
Expertise: elastin-like polypeptide gels, associative protein polymers, arrested phase separation, and neutron spin-echo measurement of relaxation in protein hydrogels.

\textbf{Darrin J. Pochan}\\
Professor of Materials Science and Engineering, University of Delaware\\
\href{mailto:pochan@udel.edu}{pochan@udel.edu}\\
Expertise: peptide and protein self-assembly, $\beta$-sheet hydrogel networks, and small-angle scattering and rheology of soft biomaterials.

\textbf{Andrei Fluerasu}\\
Lead Beamline Scientist, Coherent Hard X-ray Scattering, National Synchrotron Light Source II, Brookhaven National Laboratory\\
\href{mailto:fluerasu@bnl.gov}{fluerasu@bnl.gov}\\
Expertise: X-ray photon correlation spectroscopy and the microscopic dynamics of gels and arrested soft matter.

\textbf{Theyencheri Narayanan}\\
Senior Scientist, European Synchrotron Radiation Facility, Grenoble\\
\href{mailto:narayan@esrf.fr}{narayan@esrf.fr}\\
Expertise: small-angle X-ray scattering and X-ray photon correlation spectroscopy of colloids and soft matter.

\end{document}
